\documentclass[10pt,a4paper]{amsart}
\usepackage[margin=19mm]{geometry}
\usepackage{amsmath,amssymb,mathtools}
\usepackage[scr=rsfs,cal=euler]{mathalfa}
\usepackage{xfrac}
\usepackage[unicode,hidelinks,bookmarks=false]{hyperref}
\newcommand{\ie}{i.e.\ }
\newcommand{\cf}{cf.\ }
\newcommand{\resp}{respectively\ }
\newcommand{\Subsection}{\subsection}
\providecommand{\nobreakdash}{\nobreak\hbox{-}}
\setlength{\emergencystretch}{2em}
\multlinegap0pt
\usepackage{amssymb}
\usepackage{tikz}
\usepackage{mathtools}
\newtheorem{theorem}{Theorem}
\newcommand{\LL}{\mathcal{L}}
\newcommand{\N}{\mathbb{N}}
\newcommand{\Z}{\mathbb{Z}}
\newcommand{\eps}{\varepsilon}
\title{On the sum-of-digits measures and Cusick's conjecture via stopped random walks}
\author{Dawid Tar\l{}owski, ORCID 0000-0002-6824-4568;}
\date{Source: arXiv:2605.08624v3 (CC BY 4.0)}
\begin{document}
\maketitle

\noindent\emph{Source and licence.} On the sum-of-digits measures and Cusick's conjecture via stopped random walks. Version arXiv:2605.08624v3 (CC BY 4.0), distributed by arXiv under the Creative Commons Attribution 4.0 International licence, http://creativecommons.org/licenses/by/4.0/, as recorded in the arXiv licence metadata for this exact version and shown on the abstract page https://arxiv.org/abs/2605.08624v3. Full licence notice: \texttt{licenses/arxiv-2605.08624-CC-BY-4.0.txt}.\par\smallskip

\noindent\textit{Editorial scope.} Counted unit: the article's theorem As (source0.tex lines 843--851). The article's complete original proof of it is delivered with it (source0.tex lines 853--871, one \texttt{proof} environment of 19 lines). No mathematics is written, repaired or re-derived: the statement and the proof are the article's own lines, in the article's own notation, and no retained line is substituted or re-flowed.  No further source-local context is needed: every cross-reference of every retained line resolves inside the two delivered spans.  Nothing was omitted from any retained span, so the package carries no omission record.  No retained line contains a citation, so the wrapper carries no bibliography and no bibliography entry.  No retained line contains a figure environment, an image-inclusion command or drawing code, so no image is delivered and the package declares no asset dependency.  Editorial formatting only, in addition: an AMS \texttt{amsart} preamble that restates the article's own declarations and macros used by the retained lines, namely the environments \texttt{theorem} on separate counters, together with the standard packages those lines need.  The retained environments print in this document as Theorem 1 in order of appearance, whereas the article prints the same items with the numbering of its own full document.  The complete native source of the article as distributed in the arXiv e-print for this exact version is bundled unchanged as \texttt{sources/200163/source0.tex}.
\begin{theorem}\label{As}
  For any word $v\in\{L,R\}^{\star}$, $\lim\limits_{n\to\infty} P_{vRL^n}= \mu_v$ and $\lim\limits_{n\to\infty} P_{vLR^n}= \hat{\mu}_v$. 
Additionally: $\lim\limits_{n\to\infty}P_{L^n}=\mu_1$ and 
$\lim\limits_{n\to\infty}P_{R^n}=\hat{\mu}_1$. Hence:
 $$\{\LL(X^w_{\infty})\colon w \in\{L,R\}^\N\mbox{ is eventually constant}\}=\{\mu_t\}_{t\in T}\cup \{\hat{\mu}_t\}_{t\in T}.$$

%More precisely, for $v_1=vRL^{\infty}$ and $v_2=vLR^{\infty}$,
%$$\LL( X^{v_1}_{\infty})=\mu_{\nu}\mbox{ and }\LL(X^{v_2}_{\infty})=\hat{\mu}_{\nu}.$$
\end{theorem}

\begin{proof}
 We start with the constant sequences:  $w^1=R^{\infty}$ and $w^2=L^{\infty}$. We have 
$$T_{R^{n+1}}=[\bullet, T_R^n],\ T_{L^{n+1}}=[ T_L^n,\bullet],\mbox{ where }T_{R^0}=T_{L^0}=T_{\eps}=[\bullet,\bullet]. $$
Given $k\in\Z$, let 
$$\tau(k)=\inf\{i\in\N^+\colon \xi_i=k\}.$$ The limits of stopping times $\tau_{R^n}$, $\tau_{L^n}$, are:
$$\tau_{R^{\infty}}:=\lim\limits_{n\to\infty}\tau_{R^n}=\tau(-1),\ \tau_{L^{\infty}}:=\lim\limits_{n\to\infty}\tau_{L^n}=\tau(1).$$
 The limit $X_\infty$ is explicit: $ X_\infty^{w_1}=S_{\tau_{R^{\infty}}}=S_{\tau(-1)}\mbox{ and }X_\infty^{w_2}=S_{\tau_{L^{\infty}}}=S_{\tau(1)},$
and hence
$$\LL( X_{\infty}^{w_1})=\LL(S_{\tau(-1)})=\hat{\mu_1}\mbox{ and }\LL( X_{\infty}^{w_2})=S_{\tau(1)}=\mu_1.$$
Now, fix an arbitrary word $v\in\{L,R\}^{\star}$. We have: $T_v=[T_v^-,T_v^+]$,  $T_{vR}=[T_v^-,T_v]$, $T_{vRL}=[ T_{vR} ,T_v]$. In particular: $P_{vR}=\Phi(P_v^L,P_v)$, and:
$$P_{vRL}=\Phi(P_{vR},P_v)=\frac12\cdot \sigma_{-1}(P_{vR}) + \frac12\cdot \sigma_{1}(P_v)=\frac14\cdot \sigma_{-2}(P_v^L) + \frac14\cdot \sigma_{0}(P_{v}) + \frac12\cdot \sigma_{1}(P_{v}).$$
By iterating the above, we get:
\begin{equation}\label{iterating}
P_{vRL^n}=\frac{1}{2^{n+1}}\cdot \sigma_{-(n+1)}(P_v^L) + \sum\limits_ {i=-(n-1)}^1  \frac{1}{2^{2-i}}\cdot \sigma_{i}(P_v).
\end{equation}
  The  law  $\LL(X^w_{\infty})$, where $w=vRL^{\infty}$,  is the limit of the above sum:
$$\LL(X^w_{\infty})=\sum\limits_{i=-\infty}^{1} \frac{1}{2^{2-i}}\cdot \sigma_{i}(P_v)=\mu_1 * P_{v}=\mu_{v}.$$
To show $\LL(X_{\infty}^{w_2})=\hat{\mu_{v}}$, where $w_2=vLR^{\infty}$, we can repeat the above reasoning, or just use the symmetry $P_{\overline{w}}(d)=\hat{P}(d)$.
\end{proof}

\end{document}
